\section{Property-Relative Fork Safety}
\label{sec:fork}

A trace fork at cut $k$ retains $L_{<k}$ and gives subsequent events a child
run identity.  For a list representation,

\begin{equation}
  \fork(L,k) = \prefix(L,k).
\end{equation}

That definition makes prefix equality cheap, but it does not by itself make
re-execution or external-world claims sound.  We first describe effects, then
define the properties a fork may be asked to preserve.

\subsection{Effects along three dimensions}

A flat partition into Pure, Idempotent, Cached, and OneShot combines questions
that must be answered separately.  A cached model response describes where
replay material comes from; it does not describe whether the original call had
external cost.  We therefore use three orthogonal dimensions
(Table~\ref{tab:dimensions}).

\begin{table}[ht]
\centering
\small
\begin{tabularx}{\textwidth}{@{}>{\raggedright\arraybackslash}p{0.18\textwidth}>{\raggedright\arraybackslash}p{0.27\textwidth}Y@{}}
\toprule
Dimension & Values & Question \\
\midrule
External footprint & Pure, Idempotent, Compensatable, OneShot, Unknown & What did execution do to the external world, and can that footprint be reconciled? \\
Replay source & Deterministic, Recorded, Uncaptured & Can replay reproduce or serve the result without consulting the live oracle? \\
Lifecycle & Requested, Committed, Failed, Unknown & Is the request/outcome boundary complete and trustworthy at this cut? \\
\bottomrule
\end{tabularx}
\caption{Orthogonal effect dimensions.}
\label{tab:dimensions}
\end{table}

Pure effects are internal descriptions with no external interaction.
Idempotent effects can be repeated or overwritten under an explicit key.
Compensatable effects require a separately verified compensation.  OneShot
effects are irreversible or conservatively treated as non-repeatable.  Unknown
fails closed.  In the validation adapter, model calls are conservatively
OneShot with respect to cost and oracle interaction.  Their outputs may
simultaneously be Recorded, allowing result replay without another call.

Lifecycle is projected from events.  A request must precede exactly one
terminal outcome.  Duplicate requests, conflicting descriptors, an outcome
without a request, or a second outcome after terminal state are integrity
faults.  The first valid request and first valid terminal outcome remain the
projected authority; malformed later events do not overwrite them.

\subsection{Four fork properties}

\begin{definition}[ProjectionReplay]
The child reconstructs the same domain projection as an independently folded
copy of the retained, well-formed, ordered prefix.  No execution occurs.
\end{definition}

\begin{definition}[StrictExecutionReplay]
Every retained effect result needed by execution is deterministic or recorded,
and the cut does not retain a Requested or Unknown lifecycle.  Replaying the
prefix causes zero live oracle or tool calls.
\end{definition}

\begin{definition}[ExternalContinuation]
The runtime or caller supplies a target environment asserted to represent the
retained prefix.  The child may continue after strict replay, provided inherited
committed OneShot effects are served from the record and never executed again,
and no inherited footprint is unknown.  A snapshot identifier, fresh-environment
receipt, or equivalent attestation can discharge the environment premise; when
none is present, the executable verdict is conditional on that premise.  This
property deliberately does not assert that a retrospectively discarded source
suffix left the original world unchanged.
\end{definition}

\begin{definition}[CounterfactualWorld$(\Omega)$]
Let $\Omega$ be a declared set of external observables, such as delivered
messages, provider charges, published records, or payment-ledger entries.  In
addition to ExternalContinuation, the actual world and a world in which the
discarded suffix did not occur agree under every observation in $\Omega$.
Discarded committed OneShot or unknown effects block the claim; discarded
idempotent or compensatable effects make it conditional on reconciliation or
successful compensation.
\end{definition}

The kernel does not inspect the external world.  It checks whether retained
effect evidence is sufficient to license CounterfactualWorld$(\Omega)$ under
the stated footprint-coverage assumption.  Even a Sound verdict is therefore
a contract-relative license, not an empirical observation that two worlds are
identical.

These definitions make ``sound'' property-indexed.  The same cut may be Sound
for ProjectionReplay, Conditional for ExternalContinuation, and Unsound for
CounterfactualWorld.  Table~\ref{tab:forkconditions} summarizes the executable
assessment.

\begin{table}[ht]
\centering
\small
\begin{tabularx}{\textwidth}{@{}p{0.22\textwidth}Y Y@{}}
\toprule
Property & Retained-prefix conditions & Discarded-suffix conditions \\
\midrule
ProjectionReplay & Whole source log is structurally well formed; ordering authority is fixed & None for projected prefix equality \\
StrictExecutionReplay & No open/unknown lifecycle; every needed result is deterministic or recorded & None \\
ExternalContinuation & Strict replay; known footprint; retained committed OneShot effects are not re-executed & Target environment is attested or supplied as a premise \\
\shortstack[l]{CounterfactualWorld\\$(\Omega)$} & ExternalContinuation conditions; effect classification covers $\Omega$ & No unreconciled observed footprint; OneShot or unknown blocks \\
\bottomrule
\end{tabularx}
\caption{Property-relative fork conditions.}
\label{tab:forkconditions}
\end{table}

\subsection{Checked prefix identities}

The executable suite checks the following trace identities over generated logs.

\begin{proposition}[Identity fork]
$\fork(L,|L|)$ retains the full trace.  Under ExactTrace, the shared prefix is
equal to $L$; child lineage is stored outside that shared list.
\end{proposition}

\begin{proposition}[Nested prefix collapse]
For $j\leq i\leq |L|$, retaining $i$ events and then $j$ events has the same
shared prefix as retaining $j$ events directly.  The nested branch still records
the outer child as its parent, so raw branch records are not identical.
\end{proposition}

\begin{proposition}[Projection commutation]
For a well-formed ordered log,
\begin{equation}
 \project(\fork(L,k)) = \project(\prefix(L,k)).
\end{equation}
This is a trace/projection statement, not an external-world theorem.
\end{proposition}

\begin{proposition}[Suffix irrelevance is property-relative]
The discarded suffix is irrelevant to ProjectionReplay of the retained prefix.
It is not generally irrelevant to CounterfactualWorld.
\end{proposition}

The last proposition corrects an over-broad formulation of cheap forking.  A
committed OneShot event in the discarded suffix is a concrete witness: the
forked log excludes it and projects correctly, but the action remains true in
the world.  Likewise, a discarded request without a trustworthy terminal
outcome is Unsound because the external system may have executed it even though
the log never learned the result.

\subsection{Cuts through effects}

For each log, the validator constructs every prefix state in one forward scan
and every environmental suffix finding in one reverse scan.  It then assesses
all $|L|+1$ cuts.  Three boundary cases are especially important.

\begin{figure}[ht]
\centering
\small
\begin{tabularx}{0.96\textwidth}{@{}P{0.20\textwidth}P{0.28\textwidth}Y@{}}
\toprule
Cut & Retained prefix & Licensed claim \\
\midrule
$k_0$ before request & No inherited request & Projection is sound; a later source-side effect may still block a retrospective world claim \\
$k_1$ after request, before outcome & Requested lifecycle only & Strict replay and continuation are refused \\
$k_2$ after recorded commit & Request plus terminal recorded result & Strict replay can be sound; continuation requires zero re-execution \\
\bottomrule
\end{tabularx}
\caption{A single request/outcome pair creates three semantically different
fork regions.}
\label{fig:cut-through-request}
\end{figure}

\paragraph{Before a request.}
The retained prefix does not inherit the effect.  Projection replay is sound.
If the source later committed a OneShot effect, CounterfactualWorld is unsound
for a retrospective fork because that discarded action already occurred.

\paragraph{Between request and outcome.}
The prefix contains a Requested lifecycle.  StrictExecutionReplay and
ExternalContinuation are unsound: the effect may be running, may have committed
without acknowledgement, or may have failed.  A step bound or missing response
must not be treated as a safe cancellation.

\paragraph{After a committed OneShot.}
The response may be Recorded, making strict replay possible with zero external
calls.  ExternalContinuation is Conditional: the interpreter must serve the
recorded result and enforce zero re-executions.  This is a semantic obligation,
not merely an optimization.

These cases expose why a Decider that returns only domain events cannot express
the full agent-runtime question.  The log needs request/outcome boundaries,
footprint evidence, and replay-source evidence before it can license execution
or world-level claims.
